% diag-spcoord.tex -- diagnostico de \spcoord (TEMPORAL)
%
% Compilar:  lualatex-dev diag-spcoord.tex
% Validar:   show-pdf-tags --xml diag-spcoord.pdf | rnv-wrapp
%
% Cada caso es un parrafo numerado. Los comentarios "doc:" son lecturas
% ESPERADAS por diseno, todavia no grabadas; los casos sin "doc:" son los que
% se quiere averiguar.
%
% A  Varias coordenadas en una formula: ¿se lee la segunda? (en test-012,
%    casos 84 y 86, tres comandos con coordenada ya se leen completos).
% B  Tamano de los parentesis: \spcoord solo mira si la fraccion esta al
%    COMIENZO de la abscisa o de la ordenada. Los casos 7 a 11 hoy salen sin
%    tamano; el 6 es el control. Mirar el PDF, no la lectura.
% C  \spcoord no conoce a \spmQ: ver la lectura y el tamano.
\DocumentMetadata
  {
    lang=es-CL,
    pdfversion=2.0,
    pdfstandard=ua-2,
    tagging=on,
    tagging-setup={math/setup={mathml-SE}},
  }
\documentclass{article}
\usepackage[spanish]{babel}
\usepackage{unicode-math}
\usepackage{spintent}
\usepackage{hyperref}
\hypersetup
  {
    pdfdisplaydoctitle = true,
    pdftitle           = {diagnostico: spcoord},
  }
\pagestyle{empty}
\begin{document}

\section*{A. Varias coordenadas en una fórmula}

% doc: «1 coma 2 es igual a 3 coma 4»
1. Dos coordenadas con igual: $\spcoord{1,2} = \spcoord{3,4}$.

% doc: «1 coma 2 más 3 coma 4»
2. Dos coordenadas con suma: $\spcoord{1,2} + \spcoord{3,4}$.

% doc: «x es igual a 3 coma 4»
3. Una coordenada tras el igual, el control: $x = \spcoord{3,4}$.

% doc: «1 coma 2 es igual a 3 coma 4 es igual a 5 coma 6»
4. Tres coordenadas encadenadas: $\spcoord{1,2} = \spcoord{3,4} = \spcoord{5,6}$.

% doc: «punto A 1 coma 2 es igual a punto B 3 coma 4»
5. Dos puntos con coordenadas: $\spPoint{A}(1,2) = \spPoint{B}(3,4)$.

\section*{B. Tamaño de los paréntesis (mirar el PDF)}

% doc: «1 mitad coma 3»
6. Fracción al comienzo, el control que hoy ya dimensiona: $\spcoord{\frac{1}{2}, 3}$.

% doc: «1 más 1 mitad coma 3»
7. Fracción después de un término, en la abscisa: $\spcoord{1+\frac{1}{2}, 3}$.

% doc: «3 coma 1 más 1 mitad»
8. Fracción después de un término, en la ordenada: $\spcoord{3, 1+\frac{1}{2}}$.

% doc: «1 más 1 mitad coma 3»
9. Con dfrac después de un término: $\spcoord{1+\dfrac{1}{2}, 3}$.

10. Fracción dentro de una raíz: $\spcoord{\sqrt{\frac{1}{2}}, 1}$.

11. Fracción dentro de un subíndice: $\spcoord{x_{\frac{1}{2}}, 1}$.

\section*{C. Un número mixto dentro de una coordenada}

12. Mixto en la abscisa: $\spcoord{\spmQ{1}{1}{2}, 3}$.

13. Mixto en la ordenada: $\spcoord{3, \spmQ{1}{1}{2}}$.

14. Un mixto solo, el control: $\spmQ{1}{1}{2}$.

\section*{Con fórmula destacada}

% doc: «1 más 1 mitad coma 3»
15. Fracción después de un término, en display:
\[ \spcoord{1+\frac{1}{2}, 3} \].

\end{document}